\begin{abstract}
本論文は、VED$\times$IFGT枠組みの最初の応用層として、概念地形の場理論的記述であるConceptual Topography（CT）を提示する。
「概念地形」は、再帰的なトークンと場の相互作用のもとで提案される、情報ポテンシャル場$\Phi(x,\tau)$の長時間幾何に対するモデル内の呼称である。
沈殿、その思考との関係、見えない地形という描像は解釈であり、収束や現象的アクセスについて確立された結果ではない。

VED$\times$IFGT枠組みの中で、CTは三つの結合過程を中心に提案を整理する。
第一に、連続情報場$I(x,\tau)$を、VED$\times$IFGTの関係$I=f(1-C)$に対する認知的観測窓での提案的解釈に用い、有効ポテンシャル$\Phi = K * I$の幾何によって通りやすさを表現する。
第二に、離散的な言語トークン$\ell_i = \mathcal{Q}[I,\Phi,C;\,\text{constraints}]$は、$\nabla\Phi$との局所的相互作用を割り当てられた、再アドレス可能な離散単位として働く。再利用による井戸の深化は、定性的な機構仮説である。
第三に、再帰ループ$\Delta \to C \to I \to \Phi \to \JI \to \Delta$は、提案されたフィードバック構成を記述する。
沈殿の表式$\Phi \seddto \Phistar$は長時間テンプレートとして保持し、その収束と履歴依存的な解釈は未解決とする。

局所的な構造定義は保持する。概念は$\nabla\Phi(x) \approx 0$を満たす$\Phi$の局所的安定領域であり、理解はそのような領域を通過するコストの一時的な低下である。
思考は、形式的な流れ場$\JI$に関連する軌道に与える認知的解釈の提案であり、その場と同一ではない。
モデル内では、厳密な正値性$I(x,\tau) > 0$を、通常の完全な概念的閉包を排するものと解釈する。
さらなる分化は、残留する非対称性、境界構造、または新たな摂動のもとで提案されるものであり、正値性だけでは保証されない。

社会的スケールでは、CTはトークンの機能と支配的なコストを比較する。
共有地形、制度的アトラクタ、集合的な維持は社会的解釈の提案であり、個体と集合で同一のダイナミクスや、制度の説明が確立されたわけではない。

本論文はVED$\times$IFGT枠組みの変数を再定義しない。
その語彙の認知的・言語的な使用を提案するが、型を明示した継承、操作的な対応写像、経験的検証を確立するものではない。
\end{abstract}
